% =====================================================================
%  PETITION FOR WRIT OF MANDAMUS
%  United States Court of Appeals for the District of Columbia Circuit
%  In re McLaughlin & Fernandez
%  28 U.S.C. sec. 1651(a); Fed. R. App. P. 21; D.C. Cir. R. 21
%
%  PRO SE DRAFT.  Bracketed placeholders [ ... ] must be completed.
%  NOT FOR FILING WITHOUT REVIEW BY A LICENSED ATTORNEY.
% =====================================================================
\documentclass[14pt,letterpaper]{extarticle}

\usepackage[margin=1in]{geometry}
\usepackage{mathptmx}
\usepackage{setspace}
\usepackage{enumitem}
\usepackage{fancyhdr}
\usepackage{booktabs}
\usepackage{longtable}
\usepackage{tabularx}
\usepackage{array}
\usepackage{parskip}
\usepackage{microtype}
\usepackage{xcolor}
\usepackage[hidelinks]{hyperref}
\usepackage{ragged2e}
\usepackage{lastpage}

\renewcommand{\thesection}{\Roman{section}}
\renewcommand{\thesubsection}{\Alph{subsection}}
\setcounter{secnumdepth}{2}
\setcounter{tocdepth}{2}
\doublespacing

\pagestyle{fancy}
\fancyhf{}
\fancyhead[L]{\footnotesize In re McLaughlin \& Fernandez, No.\ \underline{\hspace{0.5in}}}
\fancyhead[R]{\footnotesize Pet.\ for Writ of Mandamus}
\fancyfoot[C]{\footnotesize Page \thepage\ of \pageref{LastPage}}
\renewcommand{\headrulewidth}{0.4pt}

\newcounter{para}
\newcommand{\np}{%
  \par\medskip
  \refstepcounter{para}%
  \noindent\hangindent=0.45in\hangafter=1
  \textbf{\arabic{para}.}\hspace{0.15in}\ignorespaces}

\newcommand{\CJP}{Md.\ Code Ann., Cts.\ \&\ Jud.\ Proc.}
\newcommand{\R}{\S}

\begin{document}
\singlespacing

% =====================================================================
\thispagestyle{empty}
\begin{center}
\fbox{\parbox{0.92\textwidth}{\centering\textbf{DRAFT --- NOT FOR FILING\\[4pt]
\small This is an unverified drafting instrument. Every bracketed item in the
form [ ... ] must be completed, and the format, length, and docketing
requirements of Fed.\ R.\ App.\ P.\ 21 and D.\ C.\ Cir.\ R.\ 21 must be
confirmed before any filing. This document was prepared with automated drafting
assistance and is not legal advice.}}}
\end{center}
\vspace{0.35in}

% =====================================================================
%  CAPTION
% =====================================================================
\begin{center}
\textbf{\large IN THE UNITED STATES COURT OF APPEALS\\[4pt]
FOR THE DISTRICT OF COLUMBIA CIRCUIT}
\end{center}
\vspace{2pt}\hrule\vspace{8pt}

\noindent
\begin{minipage}[t]{0.55\textwidth}
\raggedright
In re\\[4pt]
CAUSTIN LEE McLAUGHLIN\\
1000 North Woodington Road\\
Baltimore, Maryland 21229\\[5pt]
\hspace*{0.25in}\textit{and}\\[5pt]
JUAN JOSE FERNANDEZ\\
{[Street address]}\\
{[City, State, ZIP]}\\[8pt]
\hspace*{0.25in}Petitioners.
\end{minipage}\hfill
\begin{minipage}[t]{0.41\textwidth}
\raggedright
No.: \underline{\hspace{1.35in}}\\
(Original Proceeding ---\\
Petition for Writ of Mandamus)\\[12pt]
\textbf{PETITION FOR WRIT OF MANDAMUS}\\[8pt]
28 U.S.C.\ \R\ 1651(a) (All Writs Act)\\
Fed.\ R.\ App.\ P.\ 21\\
D.\ C.\ Cir.\ R.\ 21\\[8pt]
Filed Pro Se
\end{minipage}

\vspace{8pt}\hrule
\vspace{10pt}

% =====================================================================
{\singlespacing\tableofcontents}
\newpage
\doublespacing

% =====================================================================
\section*{Statement of the Case and Note on Jurisdiction}
\addcontentsline{toc}{section}{Statement of the Case and Note on Jurisdiction}

\np Petitioners Caustin Lee McLaughlin and Juan Jose Fernandez seek a writ of
mandamus under the All Writs Act, 28 U.S.C.\ \R\ 1651(a), and Fed.\ R.\ App.\ P.\
21, directing the performance of non-discretionary, ministerial acts and the
processing of completed excepted-service appointment packages.

\np \textbf{Note on jurisdiction and proper respondent.} A writ under the All
Writs Act may issue only ``in aid of'' the issuing court's jurisdiction, and the
Act does not itself enlarge that jurisdiction. \textit{Pennsylvania Bureau of
Correction v.\ United States Marshals Service}, 474 U.S.\ 34, 41 (1985). Relief
under the Act is discretionary and is ``sparingly exercised.'' S.\ Ct.\ R.\ 20;
see also \textit{McLaughlin v.\ United States} (Petition for Writ of Mandamus
and Prohibition filed Sept.\ 6, 2026).

\np Because the All Writs Act relief sought here concerns [federal personnel
processing / a federal agency action], Petitioners identify below the federal
officer or agency whose action would be the subject of the writ. To the extent
the writ would run against a purely private party (the Respondent trust
companies and the issuer identified in the related State petition), Petitioners
acknowledge that this Court's original jurisdiction over such a party is
doubtful, and Petitioners preserve that question for the Court. [Complete the
identification of the proper respondent before filing.]

% =====================================================================
\section{Relief Sought}
% =====================================================================

\np Petitioners seek a writ of mandamus, or in the alternative an order to show
cause, directing:

\begin{enumerate}[label=(\alph*),leftmargin=0.6in,itemsep=4pt]
\item the appropriate appointing authority for the [name the agency / office]
to complete the processing and final consideration of the completed Schedule A
and Schedule C appointment packages of Petitioners McLaughlin and Fernandez,
described in Appendix A;
\item any Respondent that is a proper party to perform the ministerial acts of
record-keeping, acknowledgment, accounting, and transmittal identified in
Appendix B; and
\item such other relief as is necessary to give effect to the foregoing.
\end{enumerate}

% =====================================================================
\section{Issues Presented}
% =====================================================================

\np Whether a writ of mandamus should issue to compel the performance of
non-discretionary, ministerial acts of record that have been requested in
conforming form and withheld.

\np Whether the failure to process completed Schedule A and Schedule C
excepted-service appointment packages to final consideration, where the
appointments are non-competitive and the processing is prescribed by regulation,
is a ministerial failure compellable by mandamus.

\np Whether Petitioners have a clear and indisputable right to the performance
of those acts and no other adequate means to obtain the relief sought.

% =====================================================================
\section{Facts Necessary to Understand the Issues}
% =====================================================================

\subsection{Petitioners}

\np Petitioner McLaughlin resides at 1000 North Woodington Road, Baltimore,
Maryland 21229. Petitioner Fernandez resides at [Street address, City, State,
ZIP] and is the spouse of Petitioner McLaughlin. Both proceed \textit{pro se}.

\np Petitioner McLaughlin asserts a medically and genomically documented
neurodevelopmental disability affecting executive function and auditory
processing under conditions of stress, and requests written instructions and
reasonable accommodation in these proceedings. See Appendix C; 42 U.S.C.\ \R\
12131 \textit{et seq.}; 28 C.F.R.\ \R\ 35.130(b)(7). He is, on information and
belief, a debtor-in-possession in an active Chapter 11, Subchapter V proceeding,
Case No.\ [24-2383].

\subsection{The Acts Requested}

\np On or about [date], Petitioner McLaughlin requested of Respondents The Bank
of New York Mellon; Delaware Trust Company; Centrus Energy Corp.\ (f/k/a USEC
Inc.); and Computershare Trust Company, N.A., the ministerial acts identified in
Appendix B, in their respective capacities as [state each capacity]. A copy of
the request is at Appendix B.

\np Respondents have not performed the requested acts and have not responded
within [number] days, notwithstanding [number] written follow-ups. See Appendix
D.

\subsection{The Excepted-Service Appointment Packages}

\np On or about [date], Petitioners submitted completed Schedule A and Schedule C
appointment packages for the positions described in Appendix A. Under 5 C.F.R.\
\R\ 213.3102, including \R\ 213.3102(u), and 5 C.F.R.\ \R\R\ 213.3301--213.3302,
appointments under Schedule A and Schedule C are excepted-service appointments
made without competitive examination. The packages have not been processed to
final consideration. See Appendix A; 5 U.S.C.\ \R\R\ 2301, 2302.

\subsection{The Integrated Implementation}

\np The appointment packages are one personnel component of a larger integrated
implementation described in Appendix E. Petitioners do not contend that
appointment is a matter of right; they contend only that the prescribed,
non-discretionary processing steps may be compelled.

% =====================================================================
\section{Reasons Why the Writ Should Issue}
% =====================================================================

\subsection{The All Writs Act Authorizes the Relief in Aid of Jurisdiction}

\np The All Writs Act authorizes the issuance of all writs ``necessary or
appropriate in aid of'' this Court's jurisdiction. 28 U.S.C.\ \R\ 1651(a). The
writ is applied sparingly, and Petitioners must show exceptional circumstances,
but the Act supplies the mechanism where a non-discretionary act has been
withheld.

\subsection{Petitioners Have a Clear and Indisputable Right}

\np A petitioner seeking mandamus must show a clear and indisputable right to
the relief sought. The acts requested are acts of record that must be performed
upon a given state of facts, in a prescribed manner, without regard to the
Respondent's own judgment upon the propriety of the act. That is the settled
definition of a ministerial duty, and it is the standard that governs the
parallel State petition. See \textit{Brack v.\ Wells}, 184 Md.\ 86, 90--91, 40
A.2d 319, 321 (1944); \textit{Philip Morris Inc.\ v.\ Angeletti}, 358 Md.\ 689,
712, 752 A.2d 200, 212 (2000); \textit{Falls Rd.\ Cmty.\ Ass'n v.\ Baltimore
County}, 439 Md.\ 1 (2014).

\subsection{Petitioners Have No Other Adequate Means}

\np An ordinary action at law would not produce the specific record acts
required and would not cure the delay. The harm is, by its nature,
non-compensable on the present record: an unprocessed appointment package
produces no final action to appeal. Petitioners therefore have no other adequate
means and invoke the established mandamus factors.

\subsection{The Writ Is Appropriate Under the Circumstances}

\np The writ is necessary because the relief sought is narrow, specific, and
non-discretionary, and because the failure to act is ongoing. Petitioners seek
only the completion of prescribed processing steps, not the outcome of any
discretionary decision.

\subsection{The Schedule A and Schedule C Processing}

\np The appointment regulations prescribe the steps an appointing authority must
take upon receipt of a completed package. Petitioners contend that those steps
are ministerial as to the acts of processing, acknowledgment, and final
transmittal, notwithstanding that eligibility and selection remain within the
appointing authority's lawful discretion.

% =====================================================================
\section{Prayer}
% =====================================================================

WHEREFORE, Petitioners respectfully request that this Court:

\begin{enumerate}[label=(\alph*),leftmargin=0.6in,itemsep=4pt]
\item issue a writ of mandamus directing the appropriate appointing authority to
complete the processing and final consideration of the Schedule A and Schedule C
packages described in Appendix A;
\item issue a writ of mandamus directing any proper Respondent to perform the
ministerial acts identified in Appendix B;
\item in the alternative, issue an order to show cause why the writ should not
issue; and
\item grant such other relief as is just and proper.
\end{enumerate}

\vspace{0.3in}
\noindent Respectfully submitted,

\vspace{0.4in}
\noindent\rule{3.0in}{0.4pt}\\
Caustin Lee McLaughlin\\
1000 North Woodington Road\\
Baltimore, Maryland 21229\\
Petitioner, Pro Se

\vspace{0.4in}
\noindent\rule{3.0in}{0.4pt}\\
Juan Jose Fernandez\\
{[Street address]}\\
{[City, State, ZIP]}\\
Petitioner, Pro Se

\vspace{0.25in}
\noindent Dated: \underline{\hspace{0.4in}} day of
\underline{\hspace{1.4in}}, 2026.

\newpage
% =====================================================================
\section{Certificate of Service}
% =====================================================================

I hereby certify that on this \underline{\hspace{0.4in}} day of
\underline{\hspace{1.4in}}, 2026, a copy of the foregoing Petition for Writ of
Mandamus was served upon each respondent and upon each party to the proceeding
below, if any, as required by Fed.\ R.\ App.\ P.\ 21(d), by [method], at the
following addresses:

\vspace{0.2in}
\noindent
\begin{tabularx}{\textwidth}{@{}p{2.6in}X@{}}
\toprule
Party / Respondent & Service address\\
\midrule
{[Name of federal agency / officer, official capacity]} & [address --- verify]\\
The Bank of New York Mellon & [address --- verify]\\
Delaware Trust Company & [address --- verify]\\
Centrus Energy Corp.\ (f/k/a USEC Inc.) & [address --- verify]\\
Computershare Trust Company, N.A. & [address --- verify]\\
\bottomrule
\end{tabularx}

\vspace{0.35in}
\noindent\rule{3.0in}{0.4pt}\\
Caustin Lee McLaughlin, Petitioner, Pro Se

\vspace{0.3in}
% =====================================================================
\section{Appendix Index}
% =====================================================================

\np Under Fed.\ R.\ App.\ P.\ 21(a)(2)(C), Petitioners attach copies of any order
or opinion, or parts of the record, essential to understand the matters set
forth in the petition.

\begin{longtable}{@{}p{1.0in}p{5.0in}@{}}
\toprule
Appendix & Description\\
\midrule
\endhead
A & Schedule A / Schedule C appointment packages and the hiring implementation
matrix for Petitioners McLaughlin and Fernandez.\\
B & The ministerial acts requested of each Respondent, with proof of submission.\\
C & [Clinical and genomic records supporting the asserted disability status.]\\
D & [Respondents' responses, or records of non-response.]\\
E & [Description of the Integrated Implementation.]\\
F & [Copies of any order or opinion, or record excerpts, required by Rule
21(a)(2)(C).]\\
\bottomrule
\end{longtable}

\vspace{0.2in}
\noindent\footnotesize\textbf{Drafting note (delete before filing).} Before any
filing in this Court, confirm (i) the identity and official capacity of the
proper federal respondent; (ii) whether D.\ C.\ Cir.\ R.\ 21 imposes a page or
word limit on the petition; (iii) the current filing fee or in forma pauperis
procedure under Fed.\ R.\ App.\ P.\ 24; and (iv) the number of copies required.

\end{document}
